% Bullet-only intermediate outline; author review pending.
\section{Assessment, Uncertainty, and Applicability}
\noindent\small Source labels: \href{run:doc_support/reconciliation/2026-11-02/SOURCES.md}{reconciliation source map}; \href{run:doc_support/reconciliation/2026-11-02/FIGURE_POINTERS.md}{figure pointers}.\normalsize

\subsection{Verification and independent evidence}
\begin{itemize}[leftmargin=*]
\item Check ROM conservation and numerical sensitivity; qualify CFD reference uncertainty separately from experimental calibration and validation.
\item Assess temperatures, circulation, duty, loss budgets, spatial reconstruction and cost separately, with engineering units and location/time meanings.
\item Qualify ACU flow evidence; CFD flow at VCU is an uncertain reference, not direct experimental circulation validation. [literature A Sections 10--11]
\item Temperature report: residuals in K by channel/location, channel-wise bias and RMSE, and declared time/space aggregation; report flow/duty and local heat budgets separately. Retain excluded channels with reasons and a sensitivity route. [MS01: score-policy lesson; Review A Section 10]
\item Use the same case names and configuration versions from source data through reduction, calibration, plotting and proposal claims. A figure caption should state whether each series is measured, CFD-referenced, calibrated or independently predicted.
\end{itemize}

\subsection{Uncertainty throughout}
\begin{itemize}[leftmargin=*]
\item Carry primitive calibration/channel, power/deposition, boundary, property, geometry and initial-state uncertainties with shared dependencies.
\item Derived duty/flow inherit their ingredients; do not count them again as independent observations.
\item Distinguish measurement variability, installed accuracy, model discrepancy and numerical uncertainty. Propagate through steady and transition outputs. [N02--N04; R01; literature A Sections 8, 15]
\item Use covariance-aware propagation when a local linear approximation is justified; reuse common calibration, property and input variables wherever they appear. Joint input sampling must preserve those same dependencies.\cite{jcgm2011vectorcovariance,jcgm2020measurementmodels}
\item Separate a scenario envelope, parameter interval and predictive interval; assess uncertainty against the intended output and available case count. Do not turn a few temperature comparisons into a calibrated rare-event probability. [Review A Section 8]
\end{itemize}

\subsection{Acceptance and limits}
\begin{itemize}[leftmargin=*]
\item Fill target/reference/metric/tolerance-symbol/runtime table; assign numbers only with an intended-use and evidence basis.
\item Sensor minima/maxima and model-inferred spatial extrema require different evidence; freezing/temperature margins require a sourced threshold. [N01; R01]
\item Discriminate physics, input, observation and numerical explanations for disagreement; report a simpler adequate domain or a supported measurement need when warranted.
\item Report local recirculation inference limits separately from bulk prediction and thermal-consequence evidence.
\end{itemize}

\subsection{Sensor changes and adequacy of the available information}
\begin{itemize}[leftmargin=*]
\item Retain the planning questions about sensor reliability, loss and evolving installation/material conditions as bounded assessment branches. [GP03/17/26; Review A Sections 5--8]
\begin{itemize}
\item Compare full and reduced sensor sets at matched inputs and fitting access; quantify the resulting loss of temperature/target accuracy and uncertainty.
\item Record calibration/drift, sensor movement, thermal expansion or changing losses as competing explanations only where evidence supports them. A mismatch is not automatically a failed sensor.
\end{itemize}
\item Connect a residual pattern to an evidence request before promising diagnosis or operator advice. [GP26--GP27; HC12--HC14]
\begin{itemize}
\item Separate input error, observation error, numerical error and model inadequacy; identify which existing or proposed measurement could distinguish them.
\item Keep ACU and VCU evidence separate. Similar processing or shared code does not make their configurations, calibration or validation interchangeable.
\end{itemize}
\end{itemize}
